%% Section V tables. Sources, all under research/gtsam_poc/ unless noted:
%%   INS, EKF-online, EKF+TL+NN : research/fgo_breadth_results.csv
%%   MPF+TL (best per case)     : research/mpf_breadth_fix_results.csv
%%   proposed estimator         : gtsam_poc_result_ps100_*.txt
%%   lag sweep                  : gtsam_poc_result_lagd*_*.txt
%%   ablation                   : gtsam_poc_result_ab*_m*.txt
%% Line metadata from the h5 exports in ~/magnav_data.

%% ---------------------------------------------------------------- Table I
\begin{table}[t]
\caption{Estimator settings. One configuration is used for every line and both
magnetometers; nothing below is tuned per case. The compensation prior
$\sigma_{\beta}$ is examined in Table~\ref{tab:ablation}.}
\label{tab:params}
\centering
\footnotesize
\setlength{\tabcolsep}{3pt}
\begin{tabular}{@{}llr@{}}
\toprule
Group & Parameter & Value \\
\midrule
Measurement & noise std.\ $\sqrt{R}$ & \SI{12}{\nano\tesla} \\
 & robust kernel, Huber $\delta$ & 1.345 \\
FOGM $S$ & $(\sigma_S,\tau)$ & \SI{3}{\nano\tesla}, \SI{180}{\second} \\
Initial $\mathbf{P}_0^{x}$ & pos.\ / alt.\ std. & \SI{0.1}{} / \SI{1.0}{\meter} \\
 & velocity std.\ (per axis) & \SI{1.0}{\meter\per\second} \\
Compensation & basis (App.~\ref{app:tl}) & P{+}I{+}E{+}bias ($m=19$) \\
 & prior $\mathbf{P}_0^{\beta}=\mathbf{I}\sigma_{\beta}^2$ & $\sigma_{\beta}=100$ \\
 & walk $\mathbf{Q}^{\beta}$ & $\mathbf{I}(1.0)^2\Delta t$ \\
Smoother & lag $L$ & \SI{300}{\second} \\
 & state interval $\Delta$ & \SI{1}{\second} ($K=10$) \\
 & solver & iSAM2, QR \\
 & process-noise floor & $10^{-20}\mathbf{I}$ \\
Evaluation & DRMS warm-up & \SI{600}{\second} \\
\bottomrule
\end{tabular}
\end{table}

%% --------------------------------------------------------------- Table II
\begin{table}[t]
\caption{Flight lines. Purpose FF = free flight (Flt1007), SV = survey
(Flt1003), CAL = calibration (Flt1006). Maps R = Renfrew\_395, E =
Eastern\_395. INS is the free-inertial horizontal drift over the line, computed
on the same grid as every other entry in this paper.}
\label{tab:lines}
\centering
\setlength{\tabcolsep}{4pt}
\begin{tabular}{llcccr}
\toprule
Line & Purp. & Map & alt.\ [m] & dur.\ [min] & INS [m] \\
\midrule
1003.02 & SV & E & 401 & 63.1 & 123.8 \\
1003.08 & SV & R & 400 & 72.4 & 271.9 \\
1006.08 & CAL & E & 398 & 14.0 & 197.7 \\
1007.02 & FF & E & 799 & 64.3 & 120.7 \\
1007.06 & FF & R & 402 & 87.3 & 317.5 \\
\bottomrule
\end{tabular}
\end{table}

%% -------------------------------------------------------------- Table III
\begin{table}[t]
\caption{Representative line 1007.06 (Flt1007, \SI{87}{\minute}), cold start on
the two uncompensated cabin magnetometers. Horizontal DRMS [m] after a
\SI{600}{\second} warm-up. The proposed estimator is reported twice: its causal
output, which uses no measurement taken after the epoch it reports and is the
like-for-like comparison against the three causal baselines, and its committed
output at a \SI{300}{\second} lag. Free-inertial drift is \SI{317.5}{\meter}.}
\label{tab:1007}
\centering
\begin{tabular}{lrrc}
\toprule
Estimator & Mag 4 & Mag 5 & look-ahead \\
\midrule
EKF, online TL & 46.7 & 17.8 & none \\
EKF{+}TL{+}NN~\cite{gnadt2022} & 48.8 & 18.3 & none \\
MPF{+}TL (best tuning) & 1227.0 & 75.1 & none \\
\midrule
\textbf{Proposed, causal} & \textbf{42.7} & \textbf{16.0} & none \\
\textbf{Proposed, committed} & \textbf{24.2} & \textbf{11.5} & \SI{300}{\second} \\
\bottomrule
\end{tabular}
\end{table}

%% --------------------------------------------------------------- Table IV
\begin{table}[t]
\caption{Breadth: cold-start cabin magnetometers across lines, flights and maps.
Horizontal DRMS [m] after a \SI{600}{\second} warm-up. ``div.''\ exceeds
\SI{10}{\kilo\meter}, ``err.''\ leaves the map. Bold marks the best causal
estimator in each row. The calibration line 1006.08 is listed for completeness
but not counted: at \SI{14}{\minute} the warm-up leaves too little line to score,
and its low map-information content makes every method weak.}
\label{tab:breadth}
\centering
\setlength{\tabcolsep}{3.5pt}
\begin{tabular}{llrrr rr}
\toprule
 & & EKF & EKF{+} & MPF{+} & \multicolumn{2}{c}{Proposed} \\
\cmidrule(l){6-7}
Line & Mag & online & TL{+}NN & TL & causal & \SI{300}{\second} \\
\midrule
1007.06 & 4 & 46.7 & 48.8 & 1227.0 & \textbf{42.7} & 24.2 \\
1007.06 & 5 & 17.8 & 18.3 & 75.1 & \textbf{16.0} & 11.5 \\
1007.02 & 4 & div. & 130.0 & 1311.5 & \textbf{74.3} & 28.4 \\
1007.02 & 5 & 31.6 & \textbf{30.4} & 2247.3 & 33.1 & 15.8 \\
1003.02 & 4 & div. & 101.0 & 482.6 & \textbf{58.2} & 19.6 \\
1003.02 & 5 & 28.1 & 29.5 & 40.9 & \textbf{21.9} & 12.0 \\
1003.08 & 4 & err. & \textbf{46.6} & 2898.7 & 48.9 & 25.0 \\
1003.08 & 5 & 21.1 & 20.7 & 77.1 & \textbf{19.2} & 10.5 \\
\midrule
1006.08 & 4 & div. & 1173.7 & --- & 134.8 & 99.1 \\
1006.08 & 5 & 117.5 & 105.2 & 793.7 & 80.1 & 101.0 \\
\midrule
\multicolumn{5}{l}{best causal, counted cases} & \textbf{6 / 8} & \\
\multicolumn{5}{l}{best overall, counted cases} & & \textbf{8 / 8} \\
\bottomrule
\end{tabular}
\end{table}

%% ---------------------------------------------------------------- Table V
\begin{table}[t]
\caption{Lag design. Horizontal DRMS [m] after a \SI{600}{\second} warm-up, for
the committed output at each lag; the causal output of the same runs is
unchanged to within \SI{1}{\meter} across the whole sweep and is given in the
last column at $L=\SI{300}{\second}$. Lengthening the lag buys accuracy in the
committed estimate and nothing in the causal one, which is what separates
smoothing gain from filtering gain.}
\label{tab:lag}
\centering
\setlength{\tabcolsep}{4pt}
\begin{tabular}{ll rrrrr r}
\toprule
 & & \multicolumn{5}{c}{lag $L$ [s], committed} & causal \\
\cmidrule(lr){3-7}
Line & Mag & 30 & 60 & 150 & 300 & 600 & \\
\midrule
1007.06 & 4 & 40.2 & 37.1 & 29.7 & 24.2 & 21.9 & 42.7 \\
1007.06 & 5 & 15.1 & 14.3 & 12.6 & 11.5 & 11.4 & 16.0 \\
1007.02 & 4 & 70.2 & 60.4 & 43.2 & 28.4 & 25.8 & 74.3 \\
1007.02 & 5 & 27.9 & 24.6 & 18.3 & 15.8 & 14.2 & 33.1 \\
1003.02 & 4 & 48.9 & 43.3 & 28.8 & 19.6 & 18.2 & 58.2 \\
1003.02 & 5 & 20.1 & 18.3 & 14.9 & 12.0 & 10.6 & 21.9 \\
1003.08 & 4 & 50.3 & 41.5 & 31.5 & 25.0 & 24.6 & 48.9 \\
1003.08 & 5 & 17.4 & 15.8 & 12.5 & 10.5 & 10.7 & 19.2 \\
\bottomrule
\end{tabular}
\end{table}

%% --------------------------------------------------------------- Table VI
\begin{table}[t]
\caption{Ablation on line 1007.06: each row changes exactly one element of the
reference configuration of Table~\ref{tab:params}. Committed DRMS [m] after a
\SI{600}{\second} warm-up, causal in parentheses. The estimator is insensitive
to every element at the metre level; the largest single effect is the
process-noise floor.}
\label{tab:ablation}
\centering
\setlength{\tabcolsep}{4pt}
\begin{tabular}{l l rr}
\toprule
Element & Change & Mag 4 & Mag 5 \\
\midrule
\multicolumn{2}{l}{reference ($\sigma_{\beta}=100$)} & 24.2 (42.7) & 11.5 (16.0) \\
\midrule
Comp.\ prior & $\sigma_{\beta}=10$ & 24.1 (43.9) & 11.5 (16.6) \\
 & $\sigma_{\beta}=1000$ & 23.9 (42.1) & 11.2 (15.4) \\
 & per column, \SI{100}{\nano\tesla} & 23.7 (42.7) & 11.6 (16.1) \\
Comp.\ walk & $\mathbf{Q}^{\beta}\times0.01$ & 26.3 (56.4) & 12.7 (20.0) \\
 & $\mathbf{Q}^{\beta}\times100$ & 24.1 (42.7) & 12.4 (25.8) \\
FOGM & $\sigma_S\times10$ & 23.9 (38.8) & 11.0 (16.6) \\
Robust kernel & removed & 25.4 (42.9) & 11.7 (16.5) \\
 & Huber $\delta=10$ & 25.6 (42.8) & 11.7 (16.5) \\
Meas.\ noise & $\sqrt{R}=\SI{65}{\nano\tesla}$ & 24.6 (42.9) & 12.7 (23.7) \\
Process floor & $10^{-16}\mathbf{I}$ & 33.1 (74.7) & 13.1 (19.3) \\
Relin. & every variable, every epoch & 24.1 (42.6) & 11.5 (16.0) \\
Meas.\ density & all \num{10} samples per state & 26.5 (44.2) & 11.7 (16.0) \\
\bottomrule
\end{tabular}
\end{table}
